\PuzzleChapterIntro{The Locksmith's Quintuple Code}{Forensics \textbullet\ Theorem Web \textbullet\ Five-Digit Lock}{This is a chain-reaction puzzle where the output of one mathematical sub-problem becomes the input parameter for the next. You must navigate Algebra, Graph Theory, Linear Algebra, and Combinatorics to derive the full code.}
\Lore{
The Locksmith’s door has a pin-code style key. It takes five dials—five small problems—each one stamped into bronze with a promise:
\begin{quote}\textit{Turn me correctly and I will not argue with the next.}\end{quote}

Each dial yields a single digit. Together they form the code that releases the latch.
}

\Problem
Five dials (I--V) each output one decimal digit.  Let those digits be \(d_1,d_2,d_3,d_4,d_5\).
The Vault code is the 5-digit integer \(\overline{d_1d_2d_3d_4d_5}\).

\medskip
\noindent\textbf{Dial I (Free-term collar).}
A collar of integer index \(a\ge 1\) used for \(n\ge 1\) full torque cycles stamps the difference
\[
\Delta=a^{n+1}-(a+1)^n.
\]
The sleeve shows \(\Delta=223137\).  Dial I outputs
\[
d_1=n.
\]

\medskip
\noindent\textbf{Dial II (Marriage minimum).}
Let \(N=2^{\,d_1+1}\).
At a party there are \(N\) boys and \(N\) girls.  Each boy knows at least \(m\) girls, and each girl knows at least \(m\) boys.
Define \(m_\ast(N)\) to be the \emph{smallest} integer \(m\) such that a complete pairing (a matching of all boys to distinct acquainted girls) is guaranteed for \emph{every} such acquaintance pattern.
Dial II outputs
\[
d_2=m_\ast(N).
\]

\medskip
\noindent\textbf{Dial III (The pegboard).}
On an \(N\times N\) grid with rows and columns indexed \(1,2,\dots,N\), a peg is present in cell \((i,j)\) iff \(i+j\) is even.
Let \(P\) be the number of ways to choose \(N\) pegs so that no two chosen pegs lie in the same row or the same column.
Dial III outputs
\[
d_3 \equiv P \pmod{10}.
\]

\medskip
\noindent\textbf{Dial IV (Commutator strip).}
Let \(t=d_2-1\).
Let \(X\) and \(B_0\) be \(t\times t\) real matrices, and define
\[
B_{k}=B_{k-1}X-XB_{k-1}\qquad(k\ge 1).
\]
Assume \(B_{t^2}=X\).
Dial IV outputs
\[
d_4=t^2.
\]

\medskip
\noindent\textbf{Dial V (Factorial determinant).}
Let \(M=d_4+1\).
Define a sequence \((u_n)_{n\ge 0}\) by \(u_0=u_1=u_2=1\) and, for every \(n\ge 0\),
\[
\det\!\begin{pmatrix}
u_{n+3} & u_{n+2}\\
u_{n+1} & u_n
\end{pmatrix}
=n!.
\]
Dial V outputs
\[
d_5 \equiv u_M \pmod{10}.
\]

\medskip
\VaultDirective{Determine the 5-digit code \(\overline{d_1d_2d_3d_4d_5}\).}

\SolutionPage
We compute the dials in order.

\smallskip
\textbf{Dial I.} If \(a^{n+1}-(a+1)^n=223137\) with integers \(a,n\ge1\), then \(a\) is an integer root of
\(f(x)=x^{n+1}-(x+1)^n-223137\), so \(a\mid f(0)=223138=2\cdot31\cdot59\cdot61\).
Modulo \(3\), \(223137\equiv0\), forcing \(a\equiv1\pmod3\) and \(n\) even.
Modulo \(4\), \(a\) must be odd and then \(a\equiv1\pmod4\).
Among odd divisors of \(31\cdot59\cdot61\), the congruences \(a\equiv1\pmod3\) and \(a\equiv1\pmod4\) force \(a=61\).
Then
\[
61^{n+1}-62^n=223137.
\]
Check \(n=2\): \(61^3-62^2=226981-3844=223137\).
If \(n\ge4\) even, then \(62^n\equiv0\pmod8\), so \(61^{n+1}\equiv223137\equiv1\pmod8\), but \(61\equiv5\pmod8\) and \(n+1\) odd gives \(61^{n+1}\equiv5\pmod8\), contradiction.
Hence \(n=2\), so \(d_1=2\).

\smallskip
\textbf{Dial II.} \(N=2^{d_1+1}=2^3=8\).
Sharpness: for even \(N\), split boys into \(B_1,B_2\) of size \(N/2\) and girls into \(G_1\) of size \(N/2-1\) and \(G_2\) of size \(N/2+1\). Connect every boy to every girl in \(G_1\), and connect \(B_1\) also to all girls in \(G_2\). Then every vertex has degree \(\ge N/2-1\), but \(B_2\) has neighborhood \(G_1\) of size \(N/2-1\), so Hall fails and no perfect matching exists. Hence \(m_\ast(N)\ge N/2\). Hall’s condition shows \(m_\ast(N)\le N/2\): if every vertex has degree \(\ge N/2\), then for any set \(X\) of boys, either \(|X|\le N/2\) and \(|N(X)|\ge N/2\ge|X|\), or \(|X|>N/2\) and then \(N(X)\) is all girls (otherwise some girl has degree \(<N/2\)). Thus \(m_\ast(8)=4\), i.e. \(d_2=4\).

\smallskip
\textbf{Dial III.} The even-parity cells form two \(4\times4\) blocks after permuting rows/cols by parity.
Choosing \(8\) nonattacking pegs means choosing a permutation in each block, so
\(P=(4!)^2=576\), hence \(d_3\equiv 6\pmod{10}\), i.e. \(d_3=6\).

\smallskip
\textbf{Dial IV.} \(t=d_2-1=3\), so \(t^2=9\).
Let \(V\) be the vector space of \(t\times t\) matrices (\(\dim V=t^2\)).
With \(T(Y)=YX-XY\), we have \(B_{k}=T^k(B_0)\).
The \(t^2+1\) vectors \(B_0,\dots,B_{t^2}\) are dependent; take a dependence with smallest index \(k\) having nonzero coefficient.
Applying \(T\) repeatedly shifts the dependence and forces \(B_{t^2}\) to be a linear combination of \(B_{t^2+1},B_{t^2+2},\dots\).
But \(B_{t^2}=X\) and \(B_{t^2+1}=[X,X]=0\), hence all later \(B\)’s are \(0\), so the only possibility is \(B_{t^2}=0\), i.e. \(X=0\).
Therefore Dial IV outputs \(d_4=t^2=9\).

\smallskip
\textbf{Dial V.} \(M=d_4+1=10\).
The determinant rule is \(u_{n+3}u_n-u_{n+2}u_{n+1}=n!\), which yields the closed form
\(u_n=(n-1)(n-3)\cdots\) (same-parity descending product), hence
\(u_{10}=9\cdot7\cdot5\cdot3\cdot1=945\), so \(d_5\equiv5\pmod{10}\), i.e. \(d_5=5\).

Thus the code is \(\overline{d_1d_2d_3d_4d_5}=24695\).

\FinalAnswer{24695}

\NextPage
